\subsection{COMPARISON RELATIONS: I. WEAK RELATIONS}

In what follows we denote by V a normed vector space over a valued field K, and by \(\mathcal{H}(\mathfrak{F}, V)\) the set of functions with values in V, each of which is defined on a subset of E belonging to the filter base F. The relations which we shall define between such functions have a local character relative to the filter with base F: let us clarify what we mean by this. If f and g are two functions from \(\mathcal{H}(\mathfrak{F}, V)\) , recall that the relation ``there is a set \(Z \in F\) such that f and g are defined and equal on Z'' is an equivalence relation \(R_{\infty}\) on \(\mathcal{H}(\mathfrak{F}, V)\) (Gen.~Top., I, p.~66). This being so, we shall say that a relation S involving a function f of \(\mathcal{H}(\mathfrak{F}, V)\) is of local character (along F) relative to f, if it is compatible (in f) with the equivalence relation \(R_{\infty}\) (Set Theory, II, p.~117); we know that if \(\tilde{f}\) is the germ of f along F, the equivalence class of f modulo \(R_{\infty}\) (an element of the quotient set \(\mathcal{H}_{\infty}(\mathfrak{F}, V) = \mathcal{H}(\mathfrak{F}, V)/R_{\infty}\) ), then one can derive from S, by passage to the quotient, a relation between \(\tilde{f}\) and the other arguments of S, and that, conversely, every relation of this nature defines a relation of local character relative to f.

Example. If f and g are two functions in \(\mathcal{H}(\mathfrak{F},\mathbf{R})\) , the relation ``there is an \(X \in F\) such that f and g are defined on X and \(f(t) \leqslant g(t)\) for every \(t \in X\) '' is of local character relative to f and g. We denote by \(\tilde{f} \leqslant \tilde{g}\) the relation obtained by passing to the quotient (for f and g); we remark that if \(\tilde{f} \leqslant \tilde{g}\) then there exist a function \(f_{1} \in \tilde{f}\) and a function \(g_{1} \in \tilde{g}\) , defined on all of E, such that \(f_{1}(t) \leqslant g_{1}(t)\) for all \(t \in E\) .

Remarks. 1) Let \(V_{i}\) ( \(1 \leqslant i \leqslant n\) ) be n normed vector spaces over K, and \(\varphi\) a function defined on \(V_{1} \times V_{2} \times \cdots \times V_{n}\) , with values in V; by passing to the quotient along \(R_{\infty}\) the function \(\varphi\) defines a map from

\[
\mathcal {H} _ {\infty} (\mathfrak {F}, V _ {1}) \times \dots \times \mathcal {H} _ {\infty} (\mathfrak {F}, V _ {n})
\]

into \(\mathcal{H}_{\infty}(\mathfrak{F}, V)\) , which one most often denotes by \(\varphi(\tilde{\mathbf{f}}_1, \ldots, \tilde{\mathbf{f}}_n)\) (Gen.~Top., I, p.~67). For example, taking \(\varphi\) to be the maps \((\mathbf{x}, \mathbf{y}) \mapsto \mathbf{x} + \mathbf{y}\) and \(\mathbf{x} \mapsto \mathbf{x}\lambda\) ( \(\lambda \in \mathbf{K}\) ), one thus defines, for any two germs \(\tilde{\mathbf{f}}\) , \(\tilde{\mathbf{g}}\) in \(\mathcal{H}_{\infty}(\mathfrak{F}, V)\) , the elements \(\tilde{\mathbf{f}} + \tilde{\mathbf{g}}\) and \(\tilde{\mathbf{f}}\lambda\) , and one verifies immediately that the rules of composition \((\tilde{\mathbf{f}}, \tilde{\mathbf{g}}) \mapsto \tilde{\mathbf{f}} + \tilde{\mathbf{g}}\) and \((\lambda, \tilde{\mathbf{f}}) \mapsto \tilde{\mathbf{f}}\lambda\) define on \(\mathcal{H}_{\infty}(\mathfrak{F}, V)\) a vector space structure over the field \(\mathbf{K}\) ; in this space \(\tilde{0}\) is the class formed by functions equal to \(0\) on a set in \(\mathfrak{F}\) , and \(-\tilde{\mathbf{f}}\) is the class formed by functions equal to \(-\mathbf{f}\) on a set in \(\mathfrak{F}\) . In the same way, if \(V\) is an algebra over \(K\) , one defines on \(\mathcal{H}_{\infty}(\mathfrak{F}, V)\) a second internal rule of composition \((\tilde{\mathbf{f}}, \tilde{\mathbf{g}}) \mapsto \tilde{\mathbf{f}}\tilde{\mathbf{g}}\) by taking \(\varphi(\mathbf{x}, \mathbf{y}) = \mathbf{x}\mathbf{y}\) ; together with the two preceding rules this defines an algebra structure over \(K\) on \(\mathcal{H}_{\infty}(\mathfrak{F}, V)\) ; if \(V\) has a unit element \(\mathbf{e}\) then \(\mathcal{H}_{\infty}(\mathfrak{F}, V)\) has as a unit element the class \(\tilde{\mathbf{e}}\) formed by the functions equal to \(\mathbf{e}\) on some set in \(\mathfrak{F}\) ; for \(\tilde{\mathbf{f}}\) to be invertible in \(\mathcal{H}_{\infty}(\mathfrak{F}, V)\) it is necessary and sufficient that for some \(\mathbf{f} \in \tilde{\mathbf{f}}\) there exists a \(Z \in \mathfrak{F}\) such that \(\mathbf{f}(t)\) is invertible in \(V\) for every \(t \in Z\) (in which case this condition is satisfied by every function in the class \(\tilde{\mathbf{f}}\) ).

\begin{enumerate}
\def\labelenumi{\arabic{enumi})}
\setcounter{enumi}{1}
\tightlist
\item
  With the same notation, let \(\psi\) be a map of a subset of \(\prod_{i=1}^{n} V_i\) into \(V\) ; we denote by \(\psi(\mathbf{f}_1, \mathbf{f}_2, \ldots, \mathbf{f}_n)\) the function equal to \(\psi(\mathbf{f}_1(t), \mathbf{f}_2(t), \ldots, \mathbf{f}_n(t))\) at every point \(t \in E\) where the \(\mathbf{f}_i(t)\) are defined and where the point \((\mathbf{f}_i(t))\) belongs to the set where \(\psi\) is defined \(^1\) . For example, \(\mathbf{f} + \mathbf{g}\) is the function equal to \(\mathbf{f}(t) + \mathbf{g}(t)\) at every point \(t \in E\) where \(\mathbf{f}\) and \(\mathbf{g}\) are both defined. Observe that the map \((\mathbf{f}, \mathbf{g}) \mapsto \mathbf{f} + \mathbf{g}\) is not a group law on \(\mathcal{H}(\mathfrak{F}, V)\) , since if \(\mathbf{f}\) is not defined on all of \(E\) there is no function \(\mathbf{g} \in \mathcal{H}(\mathfrak{F}, V)\) such that \(\mathbf{f} + \mathbf{g} = 0\) .
\end{enumerate}

DEFINITION 1. Given two real functions \(f, g\) belonging to \(\mathcal{H}(\mathfrak{F}, \mathbf{R})\) , which are \(\geqslant 0\) on a set in \(\mathfrak{F}\) , we say that \(f\) is dominated by \(g\) , or that \(g\) dominates \(f\) (along \(\mathfrak{F}\) ), and write \(f \preccurlyeq g\) or \(g \succcurlyeq f\) , if there are a set \(X \in \mathfrak{F}\) and a number \(k > 0\) such that \(f(t) \leqslant k g(t)\) for every \(t \in X\) (in other words, if there exists a \(k > 0\) such that \(\tilde{f} \leqslant k \tilde{g}\) )

Given two normed vector spaces \(V_{1}, V_{2}\) , and two functions \(f_{1}, f_{2}\) belonging to \(\mathcal{H}(\mathfrak{F}, V_{1})\) and \(\mathcal{H}(\mathfrak{F}, V_{2})\) respectively, one says that \(f_{1}\) is dominated by \(f_{2}\) (along \(F\) ), and writes \(f_{1} \preccurlyeq f_{2}\) or \(f_{2} \succcurlyeq f_{1}\) , if \(\|f_{1}\| \preccurlyeq \|f_{2}\|\) .

The relation \(f_{1} \preccurlyeq f_{2}\) is clearly of local character in \(f_{1}\) and \(f_{2}\) ; it is thus equivalent to the relation \(\tilde{f}_{1} \preccurlyeq \tilde{f}_{2}\) derived by passing to the quotient. When f and g are real functions one must be careful not to confuse the relations \(\tilde{f} \preccurlyeq \tilde{g}\) and \(\tilde{f} \leqslant \tilde{g}\) .

Note that for every scalar \(\lambda \neq 0\) the relation \(\mathbf{f}_1 \preccurlyeq \mathbf{f}_2\lambda\) is equivalent to \(\mathbf{f}_1 \preccurlyeq \mathbf{f}_2\) . If \(\mathbf{f}_1 \preccurlyeq \mathbf{f}_2\) there exists a set \(X \in \mathfrak{F}\) such that \(\mathbf{f}_1(x) = 0\) for every point \(x \in X\) where \(\mathbf{f}_2(x) = 0\) .

Examples. 1) The relation \(\mathbf{f} \preccurlyeq 1\) means that \(\mathbf{f}\) is bounded on a set of \(\mathfrak{F}\) .

\begin{enumerate}
\def\labelenumi{\arabic{enumi})}
\setcounter{enumi}{1}
\item
  For every function \(\mathbf{f}\) of \(\mathcal{H}(\mathfrak{F}, V)\) , and every scalar \(\lambda \neq 0\) , one has \(\mathbf{f} \preccurlyeq \mathbf{f}\lambda\) .
\item
  When \(x\) tends to \(+\infty\) one has \(\sin^2 x \preccurlyeq \sin x\) .
\item
  When \((x,y)\) tends to \((0,0)\) in \(\mathbf{R}^2\) one has
\end{enumerate}

\[
x y \preccurlyeq x ^ {2} + y ^ {2}.
\]

The following propositions are immediate consequences of def. 1:

PROPOSITION 1. If \(f, g, h\) are three functions in \(\mathcal{H}(\mathfrak{F}, \mathbf{R})\) then the relations \(f \preccurlyeq g\) and \(g \preccurlyeq h\) imply \(f \preccurlyeq h\) .

PROPOSITION 2. Let \(\mathbf{f}_1, \mathbf{f}_2\) be two functions in \(\mathcal{H}(\mathfrak{F}, V)\) and \(g\) a function in \(\mathcal{H}(\mathfrak{F}, \mathbf{R})\) . The relations \(\mathbf{f}_1 \preccurlyeq g\) and \(\mathbf{f}_2 \preccurlyeq g\) imply \(\mathbf{f}_1 + \mathbf{f}_2 \preccurlyeq g\) .

Further:

PROPOSITION 3. Let \(V_{1}, V_{2}, V\) be three normed spaces over the same valued field, and \((\mathbf{x}, \mathbf{y}) \mapsto [\mathbf{x}, \mathbf{y}]\) a bilinear map from \(V_{1} \times V_{2}\) into V. If \(f_{1}\) and \(f_{2}\) are functions in \(\mathcal{H}(\mathfrak{F}, V_{1})\) and \(\mathcal{H}(\mathfrak{F}, V_{2})\) respectively, and \(g_{1}, g_{2}\) are two functions in \(\mathcal{H}(\mathfrak{F}, \mathbf{R})\) such that \(f_{1} \preccurlyeq g_{1}\) and \(f_{2} \preccurlyeq g_{2}\) then \([f_{1}.f_{2}] \preccurlyeq g_{1}g_{2}\) .

Indeed (Gen.~Top., IX, p.~173, th. 1) there exists a number \(a > 0\) such that \(\|[\mathbf{f}_1.\mathbf{f}_2]\| \leqslant a\|\mathbf{f}_1\|\|\mathbf{f}_2\|\) .

COROLLARY. If V is a normed algebra, if \(f_{1}, f_{2}\) are two functions in \(\mathcal{H}(\mathfrak{F}, V)\) , and \(g_{1}, g_{2}\) are two functions in \(\mathcal{H}(\mathfrak{F}, \mathbf{R})\) , then the relations \(f_{1} \preccurlyeq g_{1}\) , \(f_{2} \preccurlyeq g_{2}\) imply \(f_{1}f_{2} \preccurlyeq g_{1}g_{2}\) .

The relation \(f \preccurlyeq g\) between functions in \(\mathcal{H}(\mathfrak{F}, V)\) is transitive by prop. 1; since it is reflexive the relation ``f \(\preccurlyeq g\) and \(g \preccurlyeq f\) '' is an equivalence relation on \(\mathcal{H}(\mathfrak{F}, V)\) (Set Theory, II, p.~113).

DEFINITION 2. Given two functions \(\mathbf{f}\) , \(\mathbf{g}\) of \(\mathcal{H}(\mathfrak{F}, V)\) we say that \(\mathbf{f}\) and \(\mathbf{g}\) are similar (along \(\mathfrak{F}\) ), and write \(\mathbf{f} \asymp \mathbf{g}\) , if \(\mathbf{f} \preccurlyeq \mathbf{g}\) and \(\mathbf{g} \preccurlyeq \mathbf{f}\) .

For every scalar \(\lambda \neq 0\) the relation \(\mathbf{f} \asymp \mathbf{g}\) is equivalent to \(\mathbf{f} \asymp \mathbf{g}\lambda\) . It implies the existence of a set \(X \in \mathfrak{F}\) such that the subset of \(X\) of points where \(\mathbf{f}(x) = 0\) is identical with the subset of \(X\) of points where \(\mathbf{g}(x) = 0\) .

Examples. 1) For a real function \(f \in \mathcal{H}(\mathfrak{F}, \mathbf{R})\) the relation \(f \asymp 1\) means that there are two numbers \(a > 0\) , \(b > 0\) such that \(a \leqslant |f(x)| \leqslant b\) on a set in \(\mathfrak{F}\) , or that the function \(\log |f|\) is bounded on a set in \(\mathfrak{F}\) : one then says that \(f\) is logarithmically bounded on a set in \(\mathfrak{F}\) .

\begin{enumerate}
\def\labelenumi{\arabic{enumi})}
\setcounter{enumi}{1}
\item
  Let V be a normed space over a non-discrete valued field K, and let \(\mathbf{f}(x)=\mathbf{a}_{0}x^{n}+\mathbf{a}_{1}x^{n-1}+\cdots+\mathbf{a}_{n}\) be a polynomial in the variable \(x\in K\) , with coefficients in V, such that \(a_{0}\neq0\) . For every vector \(b\neq0\) one has \(\mathbf{f}(x)\asymp\mathbf{b}x^{n}\) as \(|x|\) tends to \(+\infty\) .
\item
  We have seen that \(\sin^2 x \preccurlyeq \sin x\) as \(x\) tends to \(+\infty\) , but we do not have \(\sin^2 x \asymp \sin x\) , even though these functions vanish at the same points.
\item
  One has \(x^{2} + xy + y^{2} \asymp x^{2} + y^{2}\) when \((x, y)\) tends to \((0, 0)\) in \(\mathbf{R}^2\) , but not \(xy \asymp x^{2} + y^{2}\) .
\end{enumerate}

It follows immediately from prop. 3 of V, p.~213, that if \(f_1\) , \(f_2\) , \(g_1\) , \(g_2\) are functions in \(\mathcal{H}(\mathfrak{F},\mathrm{K})\) (K being any valued field) the relations \(f_1 \asymp g_1\) and \(f_2 \asymp g_2\) imply that \(f_1f_2 \asymp g_1g_2\) .

We remark that in contrast the relations \(f_{1} \asymp g_{1}\) and \(f_{2} \asymp g_{2}\) do not imply that \(f_{1} + f_{2} \asymp g_{1} + g_{2}\) , as is shown by the example \(f_{1}(x) = g_{1}(x) = x^{2}\) , \(f_{2}(x) = -(x^{2} + x)\) , \(g_{2}(x) = -(x^{2} - 1)\) , as the real variable \(x\) tends to \(+\infty\) .

The comparison relations \(f \preccurlyeq g\) , \(f \asymp g\) are said to be weak. We say that two functions f, g from \(\mathcal{H}(\mathfrak{F}, V)\) are weakly comparable if they satisfy one (at least) of the relations \(f \preccurlyeq g\) , \(g \preccurlyeq f\) .

Remarks. 1) Two functions in \(\mathcal{H}(\mathfrak{F},\mathbf{R})\) need not be weakly comparable, as is shown by the example of the functions 1 and \(x\sin x\) as \(x\) tends to \(+\infty\) .

\begin{enumerate}
\def\labelenumi{\arabic{enumi})}
\setcounter{enumi}{1}
\tightlist
\item
  Denote by \(\mathbf{R}_0\) the relation \(\mathbf{f} \asymp \mathbf{g}\) on \(\mathcal{H}(\mathfrak{F}, V)\) , and by \(\mathcal{H}_0(\mathfrak{F}, V)\) the quotient set \(\mathcal{H}(\mathfrak{F}, V)/\mathrm{R}_0\) ; note that the relation \(\mathrm{R}_{\infty}\) implies \(\mathbf{R}_0\) . Passing to the quotient the relation \(\mathbf{f} \preccurlyeq \mathbf{g}\) gives, by prop. 1 of V, p.~213, an order relation on \(\mathcal{H}_0(\mathfrak{F}, V)\) (Set Theory, III, p.~134); the preceding example shows that \(\mathcal{H}_0(\mathfrak{F}, V)\) is not totally ordered by this relation.
\end{enumerate}

